\documentclass[12pt, a4paper]{article}
\usepackage{../../../myconfig} % Gọi file cấu hình từ ngoài gốc vào

% ==========================================
% BẮT ĐẦU NỘI DUNG TÀI LIỆU
% ==========================================
\begin{document}

% Truyền 3 tham số: {Nhãn cam} {Chữ to chính giữa} {Chữ ở góc phải Header}
\titlebox{Chủ đề: Hàm số}{Hàm logarit, mũ}{Hàm số: Logarit, mũ}


\drawdashlines{20}
\newpage
\drawdashlines{25}
\newpage
\drawdashlines{25}
\newpage

% --- CÂU 1 ---
\questionLinesManual{5}{1}{Với \( a \) là số thực dương tùy ý, biểu thức \( a^{\frac{5}{3}} \cdot a^{\frac{1}{3}} \) là}{
    \begin{tasks}(4)
        \task \( a^5 \)
        \task \( a^{\frac{5}{9}} \)
        \task \( a^{\frac{4}{3}} \)
        \task \( a^2 \)
    \end{tasks}
}

% --- CÂU 2 ---
\questionLinesManual{5}{2}{Với \( a \) là số thực dương tùy ý, \( \sqrt{a^3} \) bằng}{
    \begin{tasks}(4)
        \task \( a^6 \)
        \task \( a^{\frac{3}{2}} \)
        \task \( a^{\frac{3}{3}} \)
        \task \( a^{\frac{1}{6}} \)
    \end{tasks}
}

% --- CÂU 3 ---
\questionLinesManual{5}{3}{Rút gọn biểu thức \( P = x^{\frac{1}{3}} \sqrt[6]{x} \) với \( x > 0 \).}{
    \begin{tasks}(4)
        \task \( P = \sqrt{x} \)
        \task \( P = x^{\frac{1}{8}} \)
        \task \( P = x^{\frac{2}{9}} \)
        \task \( P = x^2 \)
    \end{tasks}
}

% --- CÂU 4 ---
\questionLinesManual{5}{4}{Với mọi số thực \( a \) dương, \( \log_2 \dfrac{a}{2} \) bằng}{
    \begin{tasks}(4)
        \task \( \dfrac{1}{2} \log_2 a \)
        \task \( \log_2 a + 1 \)
        \task \( \log_2 a - 1 \)
        \task \( \log_2 a - 2 \)
    \end{tasks}
}

% --- CÂU 5 ---
\questionLinesManual{5}{5}{Với \( a \) là số thực dương tùy ý, \( 4\log \sqrt{a} \) bằng}{
    \begin{tasks}(4)
        \task \(-4\log a\)
        \task \(8\log a\)
        \task \(2\log a\)
        \task \(-2\log a\)
    \end{tasks}
}

% --- CÂU 6 ---
\questionLinesManual{5}{6}{Với \( a \) là số thực dương tùy ý, \( \log(100a) \) bằng}{
    \begin{tasks}(4)
        \task \(1 - \log a\)
        \task \(2 + \log a\)
        \task \(2 - \log a\)
        \task \(1 + \log a\)
    \end{tasks}
}

% --- CÂU 7 ---
\questionLinesManual{5}{7}{Với \( a, b \) là các số thực dương tùy ý và \( a \neq 1 \), \( \log_a \left( \sqrt[3]{b} \right) \) bằng}{
    \begin{tasks}(4)
        \task \( 3 + \log_a b \)
        \task \( 3 \log_a b \)
        \task \( \dfrac{1}{3} + \log_a b \)
        \task \( \dfrac{1}{3} \log_a b \)
    \end{tasks}
}

% --- CÂU 8 ---
\questionLinesManual{5}{8}{Tập xác định của hàm số \( y = \log_3(x - 4) \) là}{
    \begin{tasks}(4)
        \task \( (5; +\infty) \)
        \task \( (-\infty; +\infty) \)
        \task \( (4; +\infty) \)
        \task \( (-\infty; 4) \)
    \end{tasks}
}

% --- CÂU 9 ---
\questionLinesManual{5}{9}{Tập xác định của hàm số \( y = 5^x \) là}{
    \begin{tasks}(4)
        \task \( \mathbb{R} \)
        \task \( (0; +\infty) \)
        \task \( \mathbb{R} \setminus \{0\} \)
        \task \( [0; +\infty) \)
    \end{tasks}
}

% --- CÂU 10---
\questionLinesManual{5}{10}{Tìm tập xác định \( D \) của hàm số \( y = \log_2 (x^2 - 2x - 3) \).}{
    \begin{tasks}(2)
        \task \( D = (-\infty; -1] \cup [3; +\infty) \)
        \task \( D = [-1; 3] \)
        \task \( D = (-\infty; -1) \cup (3; +\infty) \)
        \task \( D = (-1; 3) \)
    \end{tasks}
}

% --- CÂU 11 ---
\questionLinesManual{5}{11}{Nghiệm của phương trình \( \log_3(5x) = 2 \) là}{
    \begin{tasks}(4)
        \task \( x = \dfrac{8}{5} \)
        \task \( x = 9 \)
        \task \( x = \dfrac{9}{5} \)
        \task \( x = 8 \)
    \end{tasks}
}

% --- CÂU 12 ---
\questionLinesManual{5}{12}{Nghiệm của phương trình \( \log_2(x+9) = 5 \) là}{
    \begin{tasks}(4)
        \task \( x = 41 \)
        \task \( x = 23 \)
        \task \( x = 1 \)
        \task \( x = 16 \)
    \end{tasks}
}

% --- CÂU 13 ---
\questionLinesManual{5}{13}{Nghiệm của phương trình \( 5^{2x-4} = 25 \) là}{
    \begin{tasks}(4)
        \task \( x = 3 \)
        \task \( x = 2 \)
        \task \( x = 1 \)
        \task \( x = -1 \)
    \end{tasks}
}

% --- CÂU 14 ---
\questionLinesManual{5}{14}{Nghiệm của phương trình \( 3^{x+2} = 27 \) là}{
    \begin{tasks}(4)
        \task \( x = -2 \)
        \task \( x = -1 \)
        \task \( x = 2 \)
        \task \( x = 1 \)
    \end{tasks}
}

% --- CÂU 15 ---
\questionLinesManual{5}{15}{Tập nghiệm \( S \) của bất phương trình \( \log_2 (2x + 3) \geq 0 \) là}{
    \begin{tasks}(4)
        \task \( S = (-\infty; -1] \)
        \task \( S = [-1; +\infty) \)
        \task \( S = (-\infty; -1) \)
        \task \( S = (-\infty; 0] \)
    \end{tasks}
}

% --- CÂU 16 ---
\questionLinesManual{5}{16}{Tập nghiệm của bất phương trình \( \log_3 (13 - x^2) \geq 2 \) là}{
    \begin{tasks}(4)
        \task \( (-\infty; -2] \cup [2; +\infty) \)
        \task \( (-\infty; 2] \)
        \task \( (0; 2] \)
        \task \( [-2; 2] \)
    \end{tasks}
}

% --- CÂU 17 ---
\questionLinesManual{5}{17}{Tập nghiệm của bất phương trình \( 3^{x^2-13} < 27 \) là}{
    \begin{tasks}(4)
        \task \( (4; +\infty) \)
        \task \( (-4; 4) \)
        \task \( (-\infty; 4) \)
        \task \( (0; 4) \)
    \end{tasks}
}

% --- CÂU 18 ---
\questionLinesManual{5}{18}{Tập nghiệm của bất phương trình \( 2^{x^2-1} < 8 \) là}{
    \begin{tasks}(4)
        \task \( (0; 2) \)
        \task \( (-\infty; 2) \)
        \task \( (-2; 2) \)
        \task \( (2; +\infty) \)
    \end{tasks}
}

\end{document}